% Chapter 14: Foundational Principles
\chapter{Foundational Principles}
\label{ch:principles}

These four principles underlie every recommendation in this handbook.
They apply to code written by AI, code written by humans,
and code written by both together.

\section{Never Fail Silently}

Every error must be explicitly reported with recovery options.
Silent failures are the most expensive kind of bug:
they produce incorrect results that look correct,
propagate through downstream systems undetected,
and surface only when the damage is difficult to reverse.

In AI-assisted development, this principle has an additional dimension:
Claude must be configured to surface problems rather than work around them.
A CLAUDE.md that says ``if you encounter an error, report it immediately
with the error message and your analysis'' produces better outcomes than
one that says ``handle errors gracefully.''

\section{Simplicity Over Complexity}

Short functions. Flat structure. Self-documenting code.
A 20-line function with a clear name is superior to a 5-line function
with three levels of abstraction.

The corollary for AI-assisted work: simple code is easier for Claude to
reason about, test, and modify correctly. Complex abstractions increase
the probability of subtle bugs in AI-generated code because each layer
of indirection is another opportunity for misunderstanding.

\section{Immutability by Default}

Return new data structures; mark mutations explicitly.
Pure functions are easier to test, easier to parallelize,
and easier for Claude to reason about.

When mutation is necessary (performance, I/O, state management),
make it explicit: name the function \code{update\_}, use mutable types,
and document the side effects. The reader should never be surprised
by what a function modifies.

\section{Fail Fast with Diagnostics}

Stop immediately on problems with full context:
what failed, what was expected, what was received,
and what the caller can do about it.

\begin{minted}{python}
def process_data(df: pd.DataFrame, min_rows: int) -> pd.DataFrame:
    if len(df) < min_rows:
        raise ValueError(
            f"Need {min_rows} rows, got {len(df)}. "
            f"Check data source or reduce min_rows parameter."
        )
    # ... process
\end{minted}

\margintip{Error messages that include both the violation and a suggested fix reduce debugging time by 80\%.}

The error message includes: what went wrong (\code{Need N, got M}),
where to look (\code{data source}), and what to do about it
(\code{reduce min\_rows}). This is the difference between a
10-second fix and a 30-minute investigation.

\vspace{1\baselineskip}

\begin{keyinsight}
These principles are not Claude-specific---they are engineering principles
that happen to become even more important when AI is generating code.
AI amplifies both good and bad patterns. If your codebase follows
these principles, Claude's contributions will follow them too.
If your codebase is silent about errors, full of mutable state,
and deeply abstracted, Claude's contributions will be the same.
The codebase is the curriculum. Make it teach the right lessons.
\end{keyinsight}
